\documentclass[11pt]{article}
\usepackage[utf8]{inputenc}
\usepackage[T1]{fontenc}
\usepackage[spanish,es-noshorthands]{babel}
\usepackage{geometry}
\usepackage{enumitem}
\usepackage{booktabs}
\usepackage{tabularx}
\usepackage{xcolor}
\usepackage{graphicx}
\usepackage{fancyhdr}
\usepackage{hyperref}

\geometry{a4paper, margin=2.3cm}
\definecolor{darkblue}{RGB}{20,40,80}
\definecolor{prblue}{RGB}{52,152,219}
\definecolor{prorange}{RGB}{230,126,34}
\hypersetup{colorlinks=true, linkcolor=prblue, urlcolor=prblue}

\newcommand{\ev}[1]{\textcolor{prblue}{[#1]}}
\newcommand{\hip}{\textcolor{prorange}{[H]}}

% ---- Datos del equipo (completar) ----------------------------------
\newcommand{\grupo}{Grupo 1}
\newcommand{\desafio}{4}
\newcommand{\producto}{ÍndicePartidas}
\newcommand{\integrantes}{Marco Ortiz \\ Angel Benavides \\ Hector Berguño \\ Vicente Carrillo \\ Martin Fuentes \\ Sarumi Mora \\ Luis Peluk}
\newcommand{\repositorio}{https://github.com/Mharko8849/Proyecto-Innova}

\setlength{\headheight}{14pt}
\pagestyle{fancy}\fancyhf{}
\fancyhead[L]{\small\textcolor{darkblue}{\textbf{Informe PEP 1 --- \grupo\ --- \producto}}}
\fancyhead[R]{\small\textcolor{darkblue}{I+E --- USACH --- Semestre 2, 2026}}
\fancyfoot[C]{\thepage}

\begin{document}

% ---- Portada --------------------------------------------------------
\begin{titlepage}
  \centering
  \includegraphics[width=3.4cm]{Logo-Color-Usach-Web.jpg}\\[0.3cm]
  {\large Universidad de Santiago de Chile}\\
  {\large Departamento de Ingeniería Informática}\\[0.2cm]
  {\normalsize Innovación y Emprendimiento --- Semestre 2, 2026}\\[1.4cm]
  {\LARGE\bfseries\color{darkblue} Informe de la PEP 1}\\[0.4cm]
  {\large Validación del problema y del modelo de negocio}\\[1.2cm]
  {\Huge\bfseries \producto}\\[0.5cm]
  {\large \grupo\ --- Desafío: \desafio}\\[1.2cm]
  {\large\bfseries Integrantes}\\[0.3cm]
  {\large \integrantes}\\[0.8cm]
  {\normalsize Repositorio: \href{\repositorio}{Proyecto-Innova}}\\[0.2cm]
  {\normalsize Profesor: Francisco Parra O.}\\[0.6cm]
  \vfill
  {\large 12 de noviembre de 2026}
\end{titlepage}

\tableofcontents
\newpage

% =====================================================================
% Borrador PEP 1. En Overleaf, usar este archivo como documento principal.
% Mantener junto a Logo-Color-Usach-Web.jpg y unit_economics.tex.
\section*{Resumen}
\addcontentsline{toc}{section}{Resumen}

ÍndicePartidas aborda el problema de los contratistas que preparan una
cotización y, al terminar el trabajo, tienen dificultades para saber
cuánto les costó ejecutarlo. Dos casos de vidriería muestran cómo se
manifiesta el problema: un independiente reportó desviaciones de entre
20 y 30\,\% en faenas complejas, y una pyme declaró desconocer su margen
real \ev{E-01}\,\ev{E-02}.

La primera entrega propuesta conecta tres acciones: cotizar un trabajo,
registrar su costo real y comparar ambos costos. Vidriería ofrece el
punto de partida mejor documentado para probar ese recorrido. La
decisión propuesta es perseverar con un flujo privado y acotado. El
prototipo está en desarrollo y todavía queda por medir si los usuarios
registran costos repetidamente y si contratan el servicio al precio
propuesto \hip.

\section{El problema}

\subsection{Enunciado}
El contratista independiente o responsable de una pyme de vidriería
necesita cotizar trabajos a medida y saber si el precio cubrió los costos
reales. E-01 tarda entre 30 minutos y 2 horas por cotización y detecta
algunos sobrecostos durante la compra de insumos o frente a imprevistos de
obra. E-02 cotiza en sesiones de 2 a 3 horas, usa un cuaderno y detecta
errores de cubicación al pagar materiales. Las entrevistas documentan
el problema durante la preparación y ejecución de trabajos, con tiempos
y costos de casos concretos \ev{E-01}\,\ev{E-02}.

\subsection{Reencuadre}
El planteamiento inicial utilizaba la cotización para reunir información
y construir un índice de costos. Las entrevistas hicieron visible un
beneficio inmediato para quien registra los datos: conocer el resultado
económico de su propio trabajo. El primer recorrido se centra por ello
en la comparación privada. El índice sectorial queda como una
oportunidad posterior, sujeta a participación voluntaria y calidad de
los datos \ev{E-01}\,\ev{E-02}\,\ev{E-03}\,\hip.

\subsection{Por qué vale la pena resolverlo}
\begin{itemize}
  \item \textbf{Dolor.} Los sobrecostos aparecen durante el trabajo y el
  margen real puede quedar sin calcular \ev{E-01}\,\ev{E-02}.
  \item \textbf{Frecuencia.} La preparación de presupuestos forma parte
  de los trabajos descritos. Falta medir el volumen mensual y la
  frecuencia de sobrecostos por cuenta \ev{E-01}\,\ev{E-02}\,\hip.
  \item \textbf{Costo actual.} Hay tiempo invertido en preparar
  cotizaciones. E-02 evaluó software especializado y lo descartó por su
  precio. El gasto efectivo en suscripciones del segmento inicial
  sigue sin cuantificarse \ev{E-01}\,\ev{E-02}\,\hip.
  \item \textbf{Búsqueda de alternativas.} E-01 probó una planilla y
  volvió al cálculo manual. E-02 evaluó software y siguió utilizando
  un cuaderno \ev{E-01}\,\ev{E-02}.
\end{itemize}

\section{Validación con usuarios}

\subsection{Con quién hablamos}
La bitácora registra seis conversaciones y el archivo de resúmenes
contiene cinco fichas con código. El contacto de climatización figura
en la bitácora y queda pendiente incorporar su ficha al conjunto
anonimizado. La simulación de florería se utilizó para explorar un rubro
configurable. Las fechas que difieren entre archivos se indican en el
anexo.

\begin{tabularx}{\linewidth}{@{}lXll@{}}
\toprule
\textbf{Código} & \textbf{Rol y sector} & \textbf{Fecha} & \textbf{Tipo} \\
\midrule
E-01 & Contratista independiente de vidriería & 02/10/2026 & Cliente objetivo \\
E-02 & Responsable de una pyme de vidriería & 02/10/2026 & Cliente objetivo \\
E-03 & Contratista general de construcción y electricidad & Octubre de 2026 & Segundo cliente potencial \\
E-04 & Profesional de ingeniería y finanzas, con experiencia en servicios técnicos & Octubre de 2026 & Experto \\
E-05 & Constructor civil del sector edificación & Octubre de 2026 & Experto técnico \\
\bottomrule
\end{tabularx}

\subsection{Hallazgos que se repiten}
Los dos casos de vidriería describen cotización manual y ausencia de una
comparación sistemática al cierre \ev{E-01}\,\ev{E-02}. El tiempo
disponible condiciona la captura: E-02 contempla registrar en terreno si
el flujo exige pocos datos, mientras que E-05 apunta a una captura simple
a cargo del administrador de obra \ev{E-02}\,\ev{E-05}. La
confidencialidad también influye en la posibilidad de compartir datos:
se mencionaron riesgos de identificación y exposición del margen
\ev{E-01}\,\ev{E-05}.

\subsection{Lo que nos sorprendió y lo que contradijo la hipótesis inicial}
E-01 concentra la ejecución en una sola persona y considera difícil
registrar a diario. E-05 indica que, en una obra con roles separados,
el registro debería recaer en el administrador o jefe de terreno. El
piloto debe distinguir quién registra y en qué momento puede hacerlo
en cada tipo de organización \ev{E-01}\,\ev{E-05}\,\hip. E-03 pidió
referencias verificables y contactos por especialidad, lo que abre una
posible línea de producto para un segundo cliente \ev{E-03}\,\hip.

\section{El experimento}

\subsection{Tarjeta de prueba}
\begin{description}[leftmargin=!,labelwidth=4cm]
  \item[Creíamos que] Un registro móvil rápido y con pocos campos
  permitiría capturar costos reales de manera recurrente \hip.
  \item[Para verificarlo] Se diseñó un formulario con trabajo, partida,
  cantidad ejecutada, materiales, costo real, horas y observación opcional.
  \item[Mediremos] Participantes, cantidad de registros, tiempo,
  necesidad de ayuda y recurrencia.
  \item[Teníamos razón si] Al menos tres de cinco participantes hacían
  tres registros y el 80\,\% de los registros se completaba en
  30 segundos o menos.
\end{description}
Fuente del diseño y del criterio: \texttt{experimento.md}.

\subsection{Ejecución}
El prototipo está en desarrollo. El equipo realizará la prueba de
registro cuando esté disponible. El piloto queda pendiente, junto con
las mediciones de tiempo y recurrencia.

\subsection{Tarjeta de aprendizaje}
\textbf{Observamos} que el esfuerzo de registrar es una preocupación
para los entrevistados \ev{E-01}\,\ev{E-02}\,\ev{E-05}.
\textbf{Frente al criterio}, todavía faltan los datos del piloto.
\textbf{Aprendimos} que la prueba debe considerar quién registra y
cómo se inserta esa tarea en el trabajo. \textbf{Por lo tanto},
mantendremos el criterio escrito y mediremos el uso en cuanto esté
listo el prototipo.

\section{Propuesta de valor y solución}

\subsection{Perfil del cliente}
\begin{tabularx}{\linewidth}{@{}lXXX@{}}
\toprule
\textbf{Actor} & \textbf{Trabajo} & \textbf{Frustración} & \textbf{Resultado esperado} \\
\midrule
Decisor de vidriería &
Cotizar a medida y saber si el trabajo cubre sus costos \ev{E-01}\,\ev{E-02} &
Recalcular a mano y descubrir sobrecostos tarde \ev{E-01}\,\ev{E-02} &
Reutilizar componentes y conocer el resultado por trabajo \ev{E-01}\,\ev{E-02}\,\hip \\
\addlinespace
Administrador de obra &
Registrar cantidades y costos de ejecución cuando existe un rol
separado de administración \ev{E-05} &
Dificultad de añadir una tarea de registro al trabajo en terreno \ev{E-05} &
Capturar cantidades y costos mediante un flujo breve \ev{E-05}\,\hip \\
\bottomrule
\end{tabularx}

\subsection{Propuesta de valor}
Para independientes y pymes de vidriería que cotizan trabajos a medida,
ÍndicePartidas propone reutilizar componentes al presupuestar y mostrar
la diferencia entre costo previsto y costo real al terminar el trabajo
\ev{E-01}\,\ev{E-02}\,\hip. La cotización reutilizable busca reducir el
recálculo, el registro breve busca reducir el esfuerzo de captura y la
comparación entrega un dato para ajustar futuras cotizaciones \hip.
Las alzas de proveedores y contingencias en terreno siguen siendo
riesgos del trabajo \ev{E-01}.

\subsection{Alcance del MVP}
\begin{tabularx}{\linewidth}{@{}XX@{}}
\toprule
\textbf{Problema} & \textbf{Función imprescindible propuesta} \\
\midrule
Rehacer cálculos en cada cotización \ev{E-01}\,\ev{E-02} &
Preparar y aceptar una cotización a partir de una partida reutilizable \hip. \\
Falta de un registro fácil de costos en ejecución \ev{E-01}\,\ev{E-02}\,\ev{E-05} &
Registrar cantidad y costo real en un trabajo asociado a la cotización \hip. \\
Dificultad para conocer la desviación al cierre \ev{E-01}\,\ev{E-02} &
Comparar costo previsto y costo real en una vista privada \hip. \\
\bottomrule
\end{tabularx}

\textbf{Después.} Configuración avanzada de componentes, historial y
versiones ampliadas, índice agregado y directorio, según lo que muestren
las pruebas \hip.

\textbf{Nunca dentro de este alcance.} ERP, remuneraciones y facturación
electrónica, según el recorte de \texttt{problema-solucion.md}.

\subsection{Mapa de historias y prototipo}
La primera entrega recorre una partida de vidriería de principio a fin:
preparar la cotización, marcarla como aceptada, abrir el trabajo,
registrar un costo real y revisar la desviación. La maqueta de este
recorrido está en desarrollo. Los montos de demostración deben quedar
identificados como ejemplos.

\subsection{Entrevistas de solución}
El equipo cuenta con declaraciones sobre configuración modular, formato
de cotización y prueba gratuita \ev{E-01}\,\ev{E-02}. Sigue pendiente
documentar la interacción con el prototipo, los puntos de confusión y
los pasos concretos que den los usuarios después de probarlo.

\section{Modelo de negocio}

\subsection{Business Model Canvas}
\begin{tabularx}{\linewidth}{@{}lX@{}}
\toprule
\textbf{Bloque} & \textbf{Propuesta y evidencia} \\
\midrule
Segmentos & Independientes y pymes que cotizan trabajos, con vidriería como primer segmento a probar \ev{E-01}\,\ev{E-02}\,\hip. \\
Propuesta de valor & Cotizar con componentes reutilizables y comparar costos privados \ev{E-01}\,\ev{E-02}\,\hip. \\
Canales & Contacto directo y referidos como primeros canales de adquisición \hip. \\
Relaciones & Autoservicio con ayuda para configurar la primera cotización \hip. \\
Ingresos & Suscripción por cuenta. La contratación efectiva está pendiente \hip. \\
Recursos & Plataforma, catálogo y datos históricos privados \hip. \\
Actividades & Mantener el flujo de cotización y registro, resguardar datos y dar soporte \hip. \\
Socios & Proveedores de alojamiento y procesamiento de pagos \hip. \\
Costos & Operación de la plataforma, soporte, cobro, captación y desarrollo \hip. \\
\bottomrule
\end{tabularx}

\subsection{Pricing}
La métrica propuesta es el volumen mensual de cotizaciones y trabajos
gestionados \hip. Los precios propuestos en \texttt{Pricing.md} son
\$9.990 para Independiente, \$24.990 para Pyme y \$59.990 para Empresa,
mensuales antes de IVA \hip. Como referencia del costo del problema,
\texttt{Pricing.md} estima \$125.000 mensuales en tiempo de cotización
con diez cotizaciones al mes, 1,25 horas por cotización y una hora
valorada en \$10.000 \hip. El tiempo por cotización se toma como
punto medio del rango de E-01. El volumen y el valor de la hora son
supuestos \ev{E-01}\,\hip. La estimación valora el tiempo dedicado.
El ahorro alcanzable se medirá en el piloto \hip.

E-01 utiliza cálculo manual tras intentar
con planillas y E-02 descartó software especializado por su precio
\ev{E-01}\,\ev{E-02}. \texttt{competencia.md} documenta licencias,
herramientas de presupuestación y proveedores que cotizan individualmente.
Esas ofertas tienen alcances y modalidades de pago distintos.

La síntesis posterior recoge aceptación verbal y preferencia por
una prueba gratuita de 30 días \ev{E-01}\,\ev{E-02}. El criterio de
\texttt{Pricing.md} pide dos compromisos verificables entre cinco
potenciales clientes. Con la evidencia disponible, ese criterio
sigue pendiente de comprobar.

\subsection{Unit economics}
En ÍndicePartidas, una unidad es una cuenta que paga una suscripción.
El objetivo de este cálculo es estimar cuánto dinero aporta esa cuenta
y cuánto tarda en recuperar lo que costó conseguirla.

\paragraph{Cuánto deja una cuenta al mes}
Los precios propuestos en \texttt{Pricing.md} son \$9.990 para el plan
Independiente y \$24.990 para el plan Pyme, ambos antes de IVA \hip.
Si el 70\,\% de las cuentas usa el primer plan y el 30\,\% usa el segundo,
el ingreso mensual promedio por cuenta (ARPA) es de \$14.490 \hip.
Este reparto 70/30 es un supuesto deliberadamente conservador para evaluar
la viabilidad mínima con predominio del plan de menor barrera económica;
no obstante, las entrevistas de campo (\texttt{E-02}, \texttt{E-03}, \texttt{E-06})
muestran mayor dolor operativo y disposición al pago en el segmento Pyme,
por lo que si el piloto concentra más cuentas en ese plan, el ARPA aumentará
favorablemente \hip.
Estimamos costos directos de \$2.500 por cuenta Independiente y \$6.000 por
cuenta Pyme para operar el servicio, dar soporte y procesar el cobro.
Este valor desglosa comisiones de pasarela de pago ($\sim$3,2\,\% por cobro recurrente),
costo de infraestructura cloud prorrateado entre cuentas activas (\$1.150 a \$2.000
por cuenta) y soporte técnico promedio prorrateado de 6 a 18 minutos mensuales por
cuenta activa (0,1 a 0,3 horas valorizadas a \$10.000 la hora de dedicación del equipo) \hip.
Con esa misma mezcla, atender una cuenta cuesta en promedio \$3.550 al mes
\hip. La diferencia es una \textbf{contribución mensual de \$10.940 por
cuenta}, equivalente al 75,5\,\% del ingreso \hip.

\paragraph{Cuánto cuesta conseguir una cuenta}
Si gastamos \$120.000 entre captación y atención de la prueba gratuita,
y cuatro cuentas empiezan a pagar, cada cuenta cuesta \$30.000 de adquirir
(CAC) \hip. Esta cifra estima entre tres y cinco contactos gestionados en el período
de prueba, con dos horas de coordinación por contacto a \$10.000 la hora de dedicación
(\$60.000 a \$100.000), más tiempo de soporte durante los 30 días (\$20.000 a \$40.000),
utilizando \$120.000 como valor central del rango \hip. Al dividir \$30.000 por
los \$10.940 que deja cada mes, obtenemos un \textbf{plazo de recuperación de 2,74 meses}
desde que la cuenta empieza a pagar \hip.

\paragraph{Cuánto aporta antes de cancelar}
Si cada mes cancela el 10\,\% de las cuentas que ya pagan
(\textit{churn}) y esa proporción se mantiene, una cuenta permanece activa
unos 10 meses en promedio \hip. En ese tiempo aporta cerca de \$109.400
(LTV), es decir, 3,65 veces lo que cuesta conseguirla (LTV/CAC) \hip.
Con esos supuestos, una cuenta cubre su costo de adquisición. Para saber
si el negocio completo es rentable todavía hay que considerar los costos
fijos de desarrollo y administración.

\paragraph{Qué pasa si menos cuentas empiezan a pagar}
Si con los mismos \$120.000 conseguimos solo dos cuentas pagadas, el costo
por cuenta sube a \$60.000, la relación LTV/CAC baja a 1,82 y la recuperación
tarda 5,48 meses \hip. En el piloto registraremos cuánto gastamos en
captación y soporte, cuántas cuentas pasan de la prueba al pago y cuántas
cancelan cada mes. Si convertir una cuenta a pago resulta más costoso de
lo esperado, probaremos una configuración inicial guiada y mediremos si
las conversiones adicionales compensan el tiempo de soporte.


\section{Lo que no se validó}
\begin{itemize}
  \item \textbf{Uso recurrente.} Ejecutar la tarjeta de prueba cuando
  esté listo el prototipo y guardar tiempos, recurrencia y ayuda.
  \item \textbf{Pago y retención.} Ofrecer el piloto con el precio
  informado y registrar compromisos verificables, conversiones y bajas.
  \item \textbf{Costo y esfuerzo real.} Pedir ejemplos recientes y
  medir tiempo, costos directos y frecuencia de cotizaciones por cuenta.
  \item \textbf{Índice agregado y nuevos rubros.} Probar disponibilidad
  y calidad de los datos con consentimiento, y probar la configuración
  con un negocio real de otro rubro.
\end{itemize}

\section{Cómo trabajó el equipo}

\subsection{Decisiones y pivotes}
El \href{https://github.com/Mharko8849/Proyecto-Innova}{repositorio del
equipo} registra las versiones del proyecto. Las fechas siguientes
corresponden a documentación registrada en Git.
\begin{itemize}
  \item \textbf{28 de octubre.} Se documentó el plan de un primer
  recorrido privado, centrado en cotizar, registrar y comparar
  (\texttt{8b93eb2}).
  \item \textbf{3 de noviembre.} Se incorporaron el modelo de negocio,
  los precios propuestos y el guion para comprobar las hipótesis
  comerciales (\texttt{05313f5}).
  \item \textbf{9 de noviembre.} Se actualizaron los trabajos de clase
  con la síntesis complementaria sobre cotización modular, prueba
  gratuita y confiabilidad del índice (\texttt{19bbd5e}).
\end{itemize}

\subsection{Uso de agentes}
Se pidió a Codex revisar la guía, cruzar los archivos locales y preparar
un borrador en LaTeX y un escenario de economía por cuenta. Para el
borrador se contrastaron la bitácora, los resúmenes, la tarjeta de prueba,
\texttt{Pricing.md} y el historial local. El equipo confirmó que E-02
corresponde a una pyme de vidriería y que el experimento se realizará
cuando esté disponible el prototipo. Los supuestos económicos están
marcados \hip. La revisión del borrador completo por el equipo queda
pendiente.

\subsection{Reparto del trabajo}
La bitácora atribuye conversaciones a Angel Benavides, Sarumi Mora y
Marco Ortiz. El historial local registra aportes a los resúmenes de
entrevista, diseño, arquitectura, base de datos, reglas, planificación y
modelo comercial. El reparto final se completará con los aportes
confirmados por cada integrante, incluido el trabajo fuera del repositorio.

\clearpage
\appendix
\section{Bitácora de contactos}
\textbf{Fuente.} \texttt{bitacora-contactos.md}. La bitácora local
registra seis intentos, seis respuestas y seis conversaciones.
La tabla resume los contactos sin identificar personas ni empresas.
Las fechas de E-03 y E-04 requieren conciliar los registros. Para E-05
se dispone del mes.

\begin{tabularx}{\linewidth}{@{}llX@{}}
\toprule
\textbf{Registro} & \textbf{Fecha} & \textbf{Resultado del contacto} \\
\midrule
E-01 & 02/10/2026 & Entrevista a un independiente de vidriería sobre cotización manual y desviaciones. \\
E-02 & 02/10/2026 & Entrevista a una pyme de vidriería sobre presupuesto, registro y margen. \\
E-03 & Octubre de 2026 & Entrevista a un contratista general sobre referencias de costos y especialistas. \\
E-04 & Octubre de 2026 & Entrevista a un profesional de ingeniería y finanzas sobre presupuestos e históricos. \\
Climatización & 09/10/2026 & Entrevista a un contratista sobre actualización de precios y control de costos. \\
E-05 & Octubre de 2026 & Entrevista a un constructor civil sobre presupuesto, ejecución y rol de captura. \\
\bottomrule
\end{tabularx}

\section{Resúmenes de entrevista}
\textbf{Fuente.} \texttt{resumen-entrevistas.md}. Las cinco fichas
registran autorización para compartir un resumen. La conversación de
climatización requiere una ficha propia y confirmación del consentimiento.

\begin{description}
  \item[E-01.] El independiente cotiza a mano, detecta desviaciones al
  comprar materiales y considera difícil registrar diariamente cuando
  él mismo ejecuta el trabajo.
  \item[E-02.] La pyme cotiza con un cuaderno, desconoce su margen
  efectivo y contempla registrar desde el teléfono si el formulario
  es rápido y pide pocos datos.
  \item[E-03.] El contratista general consulta históricos o solicita
  cotizaciones a distintos especialistas. Le interesa un índice con
  referencias confiables y contactos por especialidad.
  \item[E-04.] El profesional destaca el valor de los históricos de
  costos, el riesgo de cotizar por debajo del costo real y la necesidad
  de validar cálculos asistidos por inteligencia artificial.
  \item[E-05.] El constructor civil explica que en un contrato a suma
  alzada los cambios se gestionan mediante partidas extraordinarias.
  Propone que el administrador de obra registre los datos y pide
  proteger los márgenes comerciales.
\end{description}

\section{Prototipo}
El prototipo está en desarrollo. Este anexo incorporará las pantallas
del primer recorrido y el registro de las pruebas de solución
cuando estén disponibles.

\section{Otros archivos del repositorio}
Las fuentes de trabajo están en \texttt{Trabajos-Clases/}. El
diseño técnico se documenta en \texttt{coding-template-main/}.
El repositorio reúne los archivos citados en el cuerpo y su historial
de cambios.

\end{document}
